\documentclass[../../main.tex]{subfiles}
\begin{document}

{\bf [解]}
題意の関数
\begin{align}
 F(x)=\dfrac{1}{x}\displaystyle\int_0^x f(t)dt \label{eq:1}
\end{align}
について考える．まず$F(x)$が単調増加であることを示す．

\cref{eq:1}の両辺を $x$ で微分して
\begin{align}
F'(x)
 &=-\frac{1}{x^2}\int_0^x f(t)dt+\frac{f(x)}{x} \\
 &=\frac{1}{x}\left(f(x)-F(x)\right) \label{eq:2}
\end{align}
である．
ここで題意より$f(x)$が単調増加だから $0\le t\le x$ の時 $f(x)\ge f(t)$ である．
よって同区間で積分して
\begin{align}
 xf(x)\ge\int_0^x f(t)dt \label{eq:3}
\end{align}
したがって $0<x$ の時\cref{eq:3}の両辺を $x$ で割って，
\begin{align}
f(x)\ge F(x) \label{eq:4}
\end{align}
である．
したがって\cref{eq:2}の右辺は$0$以上だから，$F'(x)\ge 0$である．
よって $F(x)$ は単調増加であることが示された．

次に
\begin{align}
 \begin{cases}
   3F(x) = f(x) \quad (x>0) \\
   f(1) = 1
 \end{cases} \label{eq:5}
\end{align}
を満たす$f(x)$を求めよう．
\cref{eq:2}に\cref{eq:5}を代入して
\begin{align}
 f'(x) = \dfrac{2f(x)}{x} \\
\therefore 
 xf'(x) = 2f(x)
\end{align}
である．$y=f(x)$とおくと
\begin{align}
 x\dfrac{dy}{dx} = 2y \\
 \dfrac{dx}{x} = \dfrac{dy}{2y} 
\end{align}
だから，両辺積分して積分定数を$C$とすると
\begin{align}
 \log x = \dfrac{1}{2}\log y + C \label{eq:6}
\end{align}
である．ここに\cref{eq:5}を代入して
\begin{align}
 C = 0
\end{align}
だから\cref{eq:6}に代入して
\begin{align}
 \log x = \dfrac{1}{2}\log y \\
 \therefore
  y = f(x) = x^2
\end{align}
である．確かにこの関数は$0 \le x$で単調増加であり，$F(x)=\dfrac{x^2}{3}$に対して与えられた条件\cref{eq:1}も満たすから十分である．
以上から求める関数は
\begin{align}
 f(x) = x^2
\end{align}
である．

\end{document}
